\documentclass[11pt]{article}
\usepackage[T1]{fontenc}
\usepackage[utf8]{inputenc}
\usepackage{lmodern}
\usepackage[a4paper,margin=2.4cm]{geometry}
\usepackage{microtype}
\usepackage{graphicx}
\usepackage{booktabs}
\usepackage{tabularx}
\usepackage[font=small,labelfont=bf]{caption}
\usepackage[numbers,sort&compress]{natbib}
\usepackage{url}
\usepackage{xcolor}
\usepackage[colorlinks=true,linkcolor=blue!50!black,citecolor=blue!50!black,urlcolor=blue!50!black]{hyperref}
\graphicspath{{figures/}}
\input{generated/macros.tex}
\newcommand{\code}[1]{\texttt{#1}}
\makeatletter\let\inputtab\@@input\makeatother % expandable input for table bodies

\title{Running Evo 2 7B on Google Cloud and Runpod:\\
cost, GPU memory and cross-hardware agreement of variant scores}
\author{Study repository \code{rewire-bio/evo2-cloud-operations}}
\date{October 2026}

\begin{document}
\maketitle

\begin{abstract}
\noindent Clinical informatics teams evaluating the Evo 2 DNA language model need to know what it
costs to run, what hardware it needs, and whether its outputs are stable when the hardware changes.
We ran one pinned Evo 2 7B container image, unchanged, on three single-GPU platforms: a Google Cloud
\code{a3-highgpu-1g} Spot VM (H100 80 GB), a Runpod Secure Cloud H100 80 GB pod and a Runpod L40S
48 GB pod. Each platform ran Arc Institute's forward-pass correctness test, zero-shot scoring of
3,893 BRCA1 single-nucleotide variants with 8,192 bp windows, embedding extraction for an exon
classifier, a memory probe over sequence length and a same-machine repeat. All platforms passed
the correctness test (loss \PTwoLoss). The two H100s, from different providers, produced
bit-identical variant scores and classifier probabilities. The L40S produced different scores:
median absolute difference \PFiveMedianDiff{} against the H100, Spearman correlation \PFiveSpearman,
and \PFiveRankMoves{} of 3,893 variants moved by more than 1\% of the list in rank, while the AUROC
for loss-of-function variants changed only slightly (\PFiveAurocAll{} against \PTwoAurocAll).
Every machine reproduced its own scores exactly. One forward pass fitted at 32,768 bp and ran out of memory
at 131,072 bp on both the 48 GB and the 80 GB GPU; lengths in between were not tested. Scoring cost USD \PThreePerThousand{} (Runpod H100),
\PFivePerThousand{} (Runpod L40S, which varied with the host) and \PTwoPerThousand{} (Google
Cloud H100 Spot) per 1,000 variants at list prices. Zero-shot Evo 2 scores separated BRCA1 loss-of-function variants better
than the CADD v1.3 scores supplied with the BRCA1 table (AUROC \AurocEvoAll{} against
\AurocCaddAll); the difference was concentrated in noncoding variants, and on coding variants no
difference was detected. These are operational measurements, not a clinical validation.
\end{abstract}

\section{Introduction}
Evo 2 is an open DNA language model with 7B and 40B parameter checkpoints and a context of up to
1 Mb \cite{brixi2026}. Its authors report strong zero-shot variant-effect results, particularly for
noncoding and splice variants, and publish inference code, weights and notebooks \cite{arcrepo}.
They do not describe how to provision it on commercial clouds, what it costs, or how its outputs
behave across GPU types. Users have reported that 7B embeddings differ between H100 and L40S GPUs
\cite{arc177} and that long sequences exhaust GPU memory well short of the advertised context
\cite{arc160}. For a clinical team these are not curiosities: a score that changes with the
hardware cannot be locked down by an assay validation, and a context that does not fit is not a
usable context.

This study runs one fixed workload on several platforms with the same container image, and records
correctness, agreement, memory, time and cost. It also re-implements two of Arc's published notebook
workflows (BRCA1 zero-shot scoring, which Arc demonstrates with the 1B checkpoint, and exon
classification) with the 7B checkpoints, against labels and simple baselines; no published 7B
reference value was compared. The design, platforms and budget were approved in a written
protocol before any GPU run, with four dated amendments (amendment 4 records deviations from the
original run plan, acknowledged by the study owner).

\section{Methods}
\subsection{Platforms}
Table~\ref{tab:ops} lists the platforms. All are single-GPU machines. Google Cloud's
\code{a3-highgpu-1g} is sold only as Spot or Flex-start; we used Spot. Runpod pods ran in Secure
Cloud, on demand. The Google Cloud H100 is the reference platform for agreement comparisons. AWS
\code{p5.4xlarge} and \code{g6e.2xlarge} launchers are included in the repository but were not run
(no account was available; amendment 2). Runpod chose the data centre for each pod; we did not
constrain it.

\subsection{Software and model pins}
The container is built from \code{nvcr.io/nvidia/pytorch:25.04-py3} (the base of Arc's Dockerfile),
pinned by digest, with \code{evo2==0.6.0} and \code{vtx==1.1.0}. It contains PyTorch
2.7.0 (NVIDIA build), flash-attn 2.7.3 and Transformer Engine (\code{container/pip-freeze-22970720.txt}).
Because Transformer Engine is present, the 7B checkpoints ran with their configured FP8 input
projections (\code{use\_fp8\_input\_projections: True}) on every platform, not in pure bfloat16.
The image is referenced by digest on every platform. Weights were downloaded from Hugging Face at
pinned revisions and checked against SHA-256 hashes before loading, because Arc's loader uses
\code{torch.load}, which can execute code. The exon classifier (\code{schmojo/evo2-exon-classifier})
is a two-layer network; instead of loading it with \code{trust\_remote\_code} as Arc's notebook
does, we rebuilt the network locally and loaded only the pinned \code{safetensors} weights. Study
code was copied to each machine from a committed revision; all three platforms in the recorded run
used the same revision.

\subsection{Data}
All inputs are public and were fetched from Arc's repository at a pinned commit, with SHA-256
checks: GRCh37 chromosome 17; the Findlay et al.\ BRCA1 saturation genome editing table
\cite{findlay2018} (3,893 SNVs: 2,821 functional, 249 intermediate, 823 loss of function, with
consequence, CADD v1.3 and mammalian phyloP columns, as annotated by Findlay et al.); its function
scores are freely available for nonprofit use, and this study is a nonprofit use; 122 labelled exon and non-exon positions; and
Arc's four forward-pass test sequences. No patient data were used.

\subsection{Workloads}
\textbf{W0} reimplements Arc's forward-pass test with \code{evo2\_7b}; a platform passes if the mean
loss is within 0.001 of Arc's expected 0.3477. \textbf{W1} scores 8,192 bp windows centred on each
BRCA1 SNV, as in Arc's notebook: the score is the variant window's mean log-likelihood minus the
reference window's (batch size 1). The 3,893 variants share 1,326 reference windows, so 5,219
sequences are scored. \textbf{W2} extracts final-token embeddings from layer \code{blocks.26} of
\code{evo2\_7b\_base} for forward and reverse 8,192 bp sequences at 122 positions and scores them
with the pinned classifier. \textbf{W3} runs one forward pass of \code{evo2\_7b} on synthetic
sequences of 8,192, 32,768, 131,072 and 262,144 bp and stops at the first out-of-memory error.
\textbf{W4} rescores the first 200 BRCA1 variants in a new process on the same machine.

\subsection{Metrics}
Discrimination is the AUROC for loss-of-function against functional or intermediate variants,
using the negated W1 score, with 95\% intervals from 2,000 bootstrap resamples of positions
(the three SNVs at a position are kept together). Baselines from the same table need no GPU: CADD v1.3,
mammalian phyloP (amendment 1 corrected its sign before any Evo 2 output existed) and a consequence
rule (nonsense and canonical splice, then splice region, then missense, then other). Paired
differences use the same resamples. Agreement between platforms is the maximum and median absolute
difference in per-variant scores, the Spearman correlation, and the number of variants whose rank
moves by more than 1\% of the list (39 places). Peak memory is
\code{torch.cuda.max\_memory\_allocated}, with \code{nvidia-smi} sampled every 500 ms. Cost is
billed time (from pod creation on Runpod, from \code{RUNNING} on Google Cloud, to deletion) times
the list price retrieved by API at run time, plus disk.

\section{Results}
\subsection{Correctness and agreement}
All three platforms passed W0 with mean loss \PTwoLoss. The gate has bfloat16 resolution:
per-sequence losses were identical on all GPUs, but token accuracy differed on the L40S
(0.8626 against 0.8636 in the recorded run), so W0 cannot detect differences of the size seen in W1. Each platform reproduced its own W4 scores
exactly. The Runpod H100 matched the Google Cloud H100 exactly on all 3,893 variant scores and all
122 classifier probabilities (maximum absolute difference \PThreeMaxDiff). The L40S did not
(Figure~\ref{fig:agreement}): its scores differed from the reference by a median of
\PFiveMedianDiff{} and at most \PFiveMaxDiff, with Spearman correlation \PFiveSpearman; \PFiveRankMoves{}
variants moved by more than 39 rank places, and classifier probabilities differed by up to
\PFiveExonDiff. For scale, the BRCA1 variant scores themselves are of order $10^{-4}$ to $10^{-2}$.
Two earlier runs and the independent clean reproduction, each on fresh machines on the same day,
gave BRCA1 scores and exon probabilities bit-identical to the recorded run on every platform, including L40S hosts in different data
centres with different drivers and kernels (\code{evidence/checks/verification.json}). One of
those runs, 698234c1, was rejected by the harness because the coordinator changed analysis code
while it ran; its outputs are used here only for this comparison. Among loss-of-function variants
alone, the L40S and H100 rankings correlated more strongly (Spearman 0.998), and 48 of the 50
lowest-scoring variants were shared. The L40S AUROC was higher by 0.002 (95\% interval 0.0002 to
0.004, paired grouped bootstrap), a small but detectable difference
(\code{evidence/checks/verification.json}).

\begin{figure}[htbp]
\centering
\includegraphics[width=\linewidth]{fig-agreement.png}
\caption{Per-variant BRCA1 scores on each platform against the Google Cloud H100. Same image,
inputs and code. Left: Runpod H100 (identical). Right: Runpod L40S. Dashed line: identity.}
\label{fig:agreement}
\end{figure}

\subsection{BRCA1 discrimination and baselines}
Zero-shot Evo 2 7B scores gave AUROC \AurocEvoAll{} [\AurocEvoAllLo, \AurocEvoAllHi] for all SNVs,
against \AurocCaddAll{} for CADD v1.3, \AurocPhylopAll{} for phyloP and \AurocRuleAll{} for the
consequence rule (Figure~\ref{fig:auroc}). Evo 2 exceeded CADD v1.3 by \DiffCaddAll{}
[\DiffCaddAllLo, \DiffCaddAllHi]. Within subsets, the difference was concentrated in noncoding
variants (Evo 2 \AurocEvoNoncoding{} against CADD \AurocCaddNoncoding; difference
\DiffCaddNoncoding{} [\DiffCaddNoncodingLo, \DiffCaddNoncodingHi]); on coding variants no
difference was detected (Evo 2 \AurocEvoCoding, CADD \AurocCaddCoding; difference
\DiffCaddCoding{} [\DiffCaddCodingLo, \DiffCaddCodingHi]), which does not establish that they are
equal. AUROCs for subsets do not add up to the overall AUROC, which also compares variants across
subsets. The CADD scores are version 1.3, as supplied in the Findlay et al.\ table; later CADD
releases added splicing predictors, which would be expected to help most on exactly the noncoding
and splice-region variants where Evo 2 led, and were not compared. Intervals resample positions;
neighbouring positions share almost all of their 8 kb window, so the intervals may be too narrow. The exon classifier gave AUROC \PTwoAurocExon{} on 122 positions; its training
data may include these positions, so this is a pipeline check rather than an evaluation.

\begin{figure}[htbp]
\centering
\includegraphics[width=\linewidth]{fig-auroc.png}
\caption{AUROC for BRCA1 loss-of-function variants with 95\% grouped bootstrap intervals,
reference platform. Baseline scores (CADD v1.3, mammalian phyloP) come from the Findlay et al.\ table.}
\label{fig:auroc}
\end{figure}

\subsection{Memory}
Scoring 8,192 bp windows needed \PTwoPeakGib{} GiB allocated by PyTorch on every platform. A single
32,768 bp forward pass needed \PTwoThirtyTwoGib{} GiB. At 131,072 bp the forward pass ran out of
memory on both the 80 GB H100 and the 48 GB L40S (Figure~\ref{fig:memory}); lengths between
32,768 and 131,072 bp were not tested. With this code path (one full forward pass returning all
logits), the checkpoint's 1 Mb context was not reachable on either GPU; Arc points to multi-GPU
frameworks or slower forced prompting for long sequences \cite{arc160}.

\begin{figure}[htbp]
\centering
\includegraphics[width=0.75\linewidth]{fig-memory.png}
\caption{Peak GPU memory for one forward pass against sequence length. Crosses mark the memory in
use when the next allocation failed. Dotted lines show nominal capacity (80 and 48 GB); PyTorch
reported 79.2 and 44.4 GiB usable.}
\label{fig:memory}
\end{figure}

\subsection{Throughput and cost}
\begin{table}[htbp]
\centering
\footnotesize
\setlength{\tabcolsep}{4pt}
\caption{Platforms and measured operations for the recorded run. USD per 1,000 variants assumes
the BRCA1 pattern of 5,219 scored windows per 3,893 variants and list prices at run time.}
\label{tab:ops}
\begin{tabular}{@{}l l l l r r r r r@{}}
\toprule
Platform & GPU & Location & Purchase & USD/h & Windows/s & USD/1k var. & Minutes & USD \\
\midrule
\inputtab{generated/operations-table.tex}
\bottomrule
\end{tabular}
\end{table}

The two H100s scored at \PTwoRate{} and \PThreeRate{} windows per second and the L40S at
\PFiveRate{} (Table~\ref{tab:ops}). L40S throughput depended on the host: L40S pods in Runpod's
US-TX-4 data centre scored 0.92 to 0.93 windows per second in the phased and rejected runs, against
about 0.6 on the US-MO-1 hosts used in the recorded run (run log). Its cost per variant therefore ranged from below
to above that of the Runpod H100, depending on the host it was given. The Google Cloud Spot H100
cost more per variant than the Runpod on-demand H100, because its hourly price was higher at the
same throughput; Spot capacity can also be reclaimed mid-run, which did not happen here. Setup took
a large share of this short workload: pulling the 12 GB compressed image (12.11 GB of base layers,
per the registry manifest) took 491 s on Google Cloud in the recorded run, and Google's current
Deep Learning VM image needed Docker and the NVIDIA container toolkit installed first (21 s in the
recorded run, 457 s in the phased run). The completed attempts in Table~\ref{tab:ops} cost USD
\TotalRunCost; including the failed L40S attempt (USD 0.82), the recorded run cost USD 11.38
(\code{evidence/checks/verification.json}). The whole study, including phased, rejected and
reproduction runs, cost about USD 49 in cloud compute: USD 21.56 on Runpod (from its billing API)
and about USD 27 on Google Cloud (billed VM time at the Spot list price); Google Cloud billing was
not reconciled line by line.

\subsection{Operational events}
The run log records every attempt. Earlier, phased runs on the same three platforms found and
fixed four launcher problems before the recorded run: Runpod's REST API v1 had been superseded by
v2; a Runpod account without credit was refused with HTTP 402; Google's CUDA 12.9 Deep Learning VM
image ships without Docker; and it holds the NVIDIA container toolkit at an older version, so the
current package cannot be installed alongside it. Runpod placed one H100 pod in India and the L40S
in Texas, because no data centre was specified.
In the rejected run 698234c1, both Google Cloud attempts lost GPU access for new processes inside
the running container after the main workloads had finished, probably because a systemd reload on
the host revoked the container's device permissions. For the recorded run, the launcher attached
the GPU through CDI and paused automatic updates, and every workload completed. In the recorded
run, the first L40S attempt failed when the SSH session from the coordinating laptop dropped,
because the workload ran in the foreground of that session; the second attempt completed. The run
plan itself deviated from the approved protocol (more attempts in the phased runs, and several
runs rather than one); amendment 4 lists the deviations.

\section{Discussion}
\paragraph{Hardware is part of the method.} With the same image, inputs and GPU model, Evo 2 7B
outputs reproduced exactly across two providers and across processes. Changing the GPU model
changed individual scores by about as much as the typical (median) score, reordering most of a
ranked list dominated by near-zero scores, while aggregate discrimination changed by only 0.002 in
AUROC. Reference-window log-likelihoods themselves shifted by about $10^{-4}$, which is enough to
flip the sign of small scores. The most damaging variants were more stable: 48 of the 50
lowest-scoring variants were shared between the GPUs. The 7B checkpoint uses FP8 input projections
when Transformer Engine is present, and FP8 and bf16 kernels differ between GPU generations, which
is a plausible cause; we did not isolate it. In practice a validated pipeline must fix the GPU model
or revalidate when it changes, and Arc's correctness test is not sufficient on its own: the L40S
passed it.

\paragraph{Context capacity is not memory fit.} On both GPUs tested, a single forward pass fitted at
32 kb and failed at 131 kb; extrapolating the measured memory linearly suggests that about 65 kb
might fit on an 80 GB H100 but not on a 48 GB L40S, which was not tested. Pipelines that rely on long context, such as EVEE's 65 kb windows over four million
variants (about 20,000 H100-hours) \cite{evee2026}, are a different operational class from the
8 kb scoring measured here.

\paragraph{Price and residency.} A cheaper 48 GB GPU is adequate for 8 kb scoring, but on a GPU cloud
its cost per variant depends on the host it lands on, so measure throughput before committing to
a large job. Spot pricing is not automatically cheaper than a GPU cloud's
on-demand price. GPU clouds choose a location unless told otherwise; with identifiable data the
location must be pinned, and covered by a business associate agreement where HIPAA applies. EC2,
EBS and Compute Engine are covered by their providers' BAAs \cite{awshipaa,gcphipaa}; Runpod states
HIPAA and GDPR programmes \cite{runpodcompliance}. NVIDIA's hosted Evo 2 API prohibits protected
health information and production use \cite{nvidiatrial}.

\section{Limitations}
\begin{itemize}
\item Few runs per platform; timings are receipts, not benchmarks, and throughput varied between
hosts with the same GPU model. List prices change.
\item Three platforms, two GPU models. AWS was not run. Blackwell GPUs, where a user reports wrong
losses \cite{arc157}, were not tested.
\item Only the 7B checkpoints, run with FP8 input projections. The 40B model needs two H100s or one
H200 \cite{nimprereq}.
\item The CADD baseline is v1.3, from the BRCA1 table; current CADD releases with splicing
features were not compared.
\item The BRCA1 and exon workflows are re-implementations of Arc's notebooks with the 7B
checkpoints; no published 7B reference value was compared.
\item The BRCA1 comparison is one gene from one assay; the exon task may overlap the classifier's
training data. Neither says anything about clinical validity, calibration or ancestry.
\item The cause of the L40S differences was not isolated (for example, by disabling FP8).
\item Compliance terms were taken from provider pages, not audited.
\end{itemize}

\section{Availability}
Code, protocol, amendments, run log, results and this manuscript:
\url{https://github.com/rewire-bio/evo2-cloud-operations}. Container image:
\code{ghcr.io/rewire-bio/evo2-cloud-operations} (digest in \code{configs/full.json}).

\bibliographystyle{unsrtnat}
\bibliography{references}
\end{document}
